\section{The ordinary-order alternative}
\label{sec:min-ordinary}

This section converts failure of adjoint generation at a point into a
strict asymptotic inequality for ordinary orders. Lattice counting gives
the non-strict bound; the main task is to show that its equality case
already permits a nonzero fiber value to be lifted.

Fix a closed point $x\in X$. Choose regular parameters at $x$ and a
local frame of $L$, using its powers as frames of $jL$. We order Taylor
monomials first by total degree and then lexicographically, always taking
the smallest monomial. Let
\[
 H_j=\{\operatorname{in}(s):0\ne s\in R_j\}
       \subset\Z_{\geq0}^n
\]
be the set of initial exponents. Thus
$|\operatorname{in}(s)|_1=\mult_x(s)$.

The equality between the number of initial exponents and the dimension
of the section space is the elementary counting principle in
\cite[Lemma~1.4]{lazarsfeldmustata2009} and
\cite[Proposition~2.6]{kavehkhovanskii2012}.
We include its proof for the order chosen here.

\begin{lemma}
\label{lem:min-initials}
The set $H_j$ has cardinality $N_j$, and representatives of its distinct
exponents form a basis of $R_j$. There is a constant $C>0$, independent
of $j$ and $s$, such that
\begin{equation}
\label{eq:min-multiplicity}
 \mult_x(s)\leq Cj\qquad(0\ne s\in R_j).
\end{equation}
Moreover $H_j+H_q\subset H_{j+q}$.
\end{lemma}

\begin{proof}
Taylor expansion is injective on $R_j$: a section with zero germ vanishes
on a nonempty open set and hence on the integral variety $X$. Each
monomial coefficient has a one-dimensional space of possible values.
Successively selecting a smallest initial monomial and eliminating its
coefficient from the remaining basis vectors therefore produces a basis
with distinct initials. Any linear combination of these vectors has the
initial of its first nonzero summand in this ordering. This proves the
cardinality assertion.

Choose a fixed sufficiently ample multiple $kL$. General members of its
linear system through $x$ cut out a curve smooth at $x$ and not contained
in the divisor of a prescribed nonzero section of $jL$. Local
intersection multiplicity bounds $\mult_x(s)$ from above by the degree
of that divisor on the curve, at most $jk^{n-1}L^n$. The same assertion
in dimension one is simply the degree bound for an effective divisor
on $X$. This gives \eqref{eq:min-multiplicity}. Finally, the chosen
monomial ordering is multiplicative, so initial exponents add under
products of sections.
\end{proof}

\begin{proposition}[Ordinary-order alternative]
\label{prop:ordinary-gap}
If $P$ is not generated by its global sections at $x$, then
\begin{equation}
\label{eq:min-gap}
 \liminf_{j\longrightarrow\infty}
   \frac{1}{jN_j}\sum_{\gamma\in H_j}|\gamma|_1
       >\frac{n}{n+1}.
\end{equation}
\end{proposition}

\begin{proof}
We prove that failure of the strict inequality implies generation at
$x$. The proof has three steps: lattice equality forces volume one,
the resulting leading forms give a strict threshold on the exceptional
divisor of the point blowup, and vanishing lifts a nonzero section from
that divisor. Put $V=L^n$. The number of nonnegative integral points of total
degree at most $r$ is $\binom{r+n}{n}$. Filling the points in order of
increasing total degree, and using
$N_j=Vj^n/n!+O(j^{n-1})$, gives
\begin{equation}
\label{eq:min-lattice-lower}
 \liminf_{j\longrightarrow\infty}
 \frac{1}{jN_j}\sum_{\gamma\in H_j}|\gamma|_1
 \geq V^{1/n}\frac{n}{n+1}.
\end{equation}
Indeed the last occupied degree for a set of $N_j$ smallest points is
$V^{1/n}j+O(1)$, and summing total degree over these points gives mean
degree $V^{1/n}nj/(n+1)+O(1)$. Since $V$ is a positive integer, failure
of \eqref{eq:min-gap} forces $V=1$ and equality along a subsequence.

Let
\[
 \Delta=\{u\in\R_{\geq0}^n:|u|_1\leq1\}.
\]
Along this subsequence, every fixed open ball with closure in the
interior of $\Delta$ meets $j^{-1}H_j$ for all sufficiently large $j$.
To see this, compare $H_j$ with a set $T_j$ of $N_j$ smallest lattice
points. If such a ball were missed, at least $c_1j^n$ points of $T_j$
of degrees at most $(1-c_2)j$ would be absent, for fixed positive
$c_1,c_2$. Every replacement outside $T_j$ has degree at least
$j-O(1)$. The remaining replacements cannot decrease the degree sum,
by the definition of $T_j$. The normalized mean would therefore exceed
the minimum by a fixed positive amount, contradicting equality along
the subsequence.

Choose a rational number $t'$ with $n<t'<n+1$ and put
$c=(1/t',\ldots,1/t')$. This point is in the interior of $\Delta$.
Choose finitely many small balls whose closures lie there and such that
any choice of one point from each ball has convex hull containing $c$
in its interior. For one sufficiently large degree $j$ on the
subsequence, choose initial exponents $\gamma_\nu\in H_j$ in these
balls after scaling by $j^{-1}$. The rational polytope
\[
 Q=\conv\{\gamma_\nu/j\}
\]
contains $c$ in its interior. Its intersection with
$\{|u|_1=n/t'\}$ has rational points $p_1,\ldots,p_b$ whose convex
hull contains $c$ in its interior relative to this full slice
hyperplane. In dimension one this slice is the singleton $\{c\}$.

Write each $p_\ell$ as a rational convex combination of the
$\gamma_\nu/j$, and choose sections $s_\nu$ with those initials.
If $t'=a/b'$ in lowest terms, choose a positive integer $Q_0$ divisible
by $a$ and by every denominator in these convex combinations. Taking
the corresponding products of $Q_0$ sections gives sections
$r_1,\ldots,r_b$ of $jQ_0L$. Set $h=jQ_0/t'$, which is a positive
integer by construction. Their initial exponents $\beta_\ell$ satisfy
\begin{equation}
\label{eq:min-product-initials}
 r_\ell\in H^0(X,ht'L),\qquad
 \beta_\ell=jQ_0p_\ell,\qquad |\beta_\ell|_1=nh.
\end{equation}
In particular, $\conv\{\beta_\ell/h\}$ contains
$\mathbf1=(1,\ldots,1)$ in relative interior in the full hyperplane
$\{|u|_1=n\}$.

Let $p:B=\Bl_xX\longrightarrow X$ and
$E\simeq\mathbb P^{n-1}$ be the exceptional divisor. When $n=1$, take
$p$ to be the identity and $E=x$. Removing the common order $nh$ from
the pullbacks of the $r_\ell$ gives a subsystem of
\[
 H=p^*(ht'L)-nhE
\]
with base ideal $\mathfrak b$. Its restriction to $E$ is generated,
in local frames, by the leading homogeneous forms $g_\ell$ of the
$r_\ell$, each of degree $nh$. We claim that
\begin{equation}
\label{eq:min-restricted-threshold}
 \lct_z(E,\mathfrak b|_E)>1/h\qquad(z\in E).
\end{equation}
Here the restricted ideal means its image in $\cO_E$. For $n=1$ it
is the unit ideal, so the assertion holds immediately.

The monomial threshold step below is the unit-monomial special case
of Howald's multiplier-ideal formula~\cite[Main Theorem]{howald2001};
we give the direct SNC argument in this setting.

For $n>1$, first replace the $g_\ell$ by their initial monomials
$z^{\beta_\ell}$. On the affine chart of $E$ where $z_k\ne0$, their
exponents are obtained by omitting coordinate $k$. Projection from
the hyperplane $|u|_1=n$ is an affine isomorphism, so the resulting
normalized exponent polytope contains the all-ones vector in ordinary
interior. Consequently, for every nonzero
$u\in\R_{\geq0}^{n-1}$,
\begin{equation}
\label{eq:min-monomial-strict}
 \frac1h\min_\ell\langle u,\beta_\ell^{(k)}\rangle
       <\sum_i u_i.
\end{equation}
For a prime-divisor order over this chart, use its values on the
coordinate functions as $u$. The left side is the order of the
monomial ideal divided by $h$, while the right side is at most the log
discrepancy, by the SNC inequality in Lemma~\ref{lem:retraction}. If
$u=0$, the ideal order is zero and the discrepancy is positive. Thus
the monomial ideal is klt at coefficient $1/h$. Computing its threshold
on a log resolution, where a finite minimum suffices, gives a strict
threshold greater than $1/h$ at every point.

This statement transfers to the original forms by an algebraic
degeneration. Since only finitely many monomials occur, choose an
integral weight vector $w$ which makes the lexicographic initial of
each $g_\ell$ its unique term of smallest weight. The forms
\[
 g_{\ell,\lambda}(z)
   =\lambda^{-\langle w,\beta_\ell\rangle}
      g_\ell(\lambda^{w_1}z_1,\ldots,\lambda^{w_n}z_n)
\]
have polynomial coefficients in $\lambda$, with monomial special
fiber. They define an algebraic family of ideals on
$E\times\mathbb A^1$. By Lemma~\ref{lem:lct-family}, the locus where
the fiber threshold is at most $1/h$ is closed. Its image in
$\mathbb A^1$ is closed because $E$ is proper, and does not contain
$0$. For some nonzero parameter the threshold is therefore greater
than $1/h$ everywhere on $E$. All nonzero fibers are obtained from the
original forms by a diagonal automorphism of $E$ and nonzero scalar
changes of generators. This proves
\eqref{eq:min-restricted-threshold}. This application uses precisely
the family semicontinuity of complex singularity exponents
\cite{dk2001} incorporated in Lemma~\ref{lem:lct-family}.

It remains to lift from $E$. The blowup formula gives
\begin{equation}
\label{eq:min-adjoint-identity}
 p^*P-(K_B+E)-H/h=(n+1-t')p^*L.
\end{equation}
Take a projective log resolution
$\sigma:T\longrightarrow B$ of $E$ and $\mathfrak b$, with strict
transform $E'$ and
$\mathfrak b\cO_T=\cO_T(-G)$. The divisor $E'$ is not a component
of $G$, since the restricted system is not identically zero. Set
\[
 D_T=K_T+\sigma^*(p^*P-K_B-E)-\lfloor G/h\rfloor.
\]
This is an integral divisor, and
\begin{equation}
\label{eq:min-preliminary-positivity}
 D_T-K_T-\{G/h\}
 =(n+1-t')(p\sigma)^*L+(\sigma^*H-G)/h.
\end{equation}
The first summand is nef and big and the second is nef, because
$\sigma^*H-G$ is globally generated. The fractional boundary has SNC
support. Theorem~\ref{thm:kv}, in the stated logarithmic form of
Kawamata--Viehweg vanishing \cite{fujinoKV}, therefore gives
$H^1(T,\cO_T(D_T))=0$.

Let $\sigma_E:E'\longrightarrow E$ be the induced morphism. Adjunction
identifies the restriction of $D_T+E'$ with
\begin{equation}
\label{eq:min-preliminary-restriction}
 \sigma_E^*(p^*P|_E)
   +K_{E'/E}-\lfloor G|_{E'}/h\rfloor.
\end{equation}
Indeed $E'+\Supp(G)$ is SNC. Each component restricts with
multiplicity one, and two distinct components cannot restrict to the
same prime divisor on $E'$, so taking floors commutes with restriction.
The generators on $E'$ generate $\cO_{E'}(-G|_{E'})$, the pullback
of the restricted ideal. The strict threshold in
\eqref{eq:min-restricted-threshold} implies
\[
 K_{E'/E}-\lfloor G|_{E'}/h\rfloor\geq0.
\]
The first line bundle in \eqref{eq:min-preliminary-restriction} is
trivial with fiber $P_x$. A nonzero vector in this fiber, multiplied
by the canonical section of the effective correction, gives a section
on $E'$ nonzero at its generic point. The exact sequence for $E'$ and
the vanishing just proved lift it to a section of $D_T+E'$.

Relative to $\sigma^*p^*P$, the correction of this lifted bundle is
\[
 C_T=K_{T/B}-\sigma^*E+E'-\lfloor G/h\rfloor.
\]
Its coefficient on $E'$ is zero, and its coefficients on every other
$\sigma$-nonexceptional divisor are nonpositive. Its positive part is
therefore exceptional over $X$. Since $E'$ maps to $x$ and the lifted
section is nonzero at its generic point, Lemma~\ref{lem:descent}, applied
to $p\sigma$, gives a section of $P$ nonzero at $x$. This proves the
proposition, including dimension one, when $E'$ is the point itself.
\end{proof}

\section{Finite-degree exponential minimization}
\label{sec:min-exponential}

We now allow both the basis and the monomial test to vary. The strict
ordinary-order inequality makes the objective less than one, and a
supply of low-order jets bounds the discrepancy on that sublevel set.
To prove attainment, we first select one flag and its adapted basis;
only then do we choose the resolution on which compactness is used.

For $N_m>0$, an ordered basis $\mathbf s=(s_1,\ldots,s_{N_m})$ of
$R_m$, a nonzero monomial test $v$, and $t>0$, put
\begin{equation}
\label{eq:min-objective}
 F_{m,t}(v,\mathbf s)
  =\frac1{N_m}\sum_{i=1}^{N_m}
       \exp\!\left(\frac{\A(v)}t-\frac{v(s_i)}m\right),
 \qquad
 f_{m,t}(z)=\inf_{\substack{v,\mathbf s\\z\in c_X(v)}}
                F_{m,t}(v,\mathbf s).
\end{equation}
Here $c_X(v)$ denotes the closed center, and the infimum also ranges
over the bases. Positive rescaling of a test is always allowed.
Approximating its positive weights by rational weights preserves its
center and approximates the orders of finitely many sections and its
discrepancy. Thus the same infimum is obtained using positive multiples
of prime-divisor orders.

\begin{lemma}
\label{lem:min-order-bound}
If the center of a test $v$ contains $x$, then every nonzero section
$s$ satisfies
\begin{equation}
\label{eq:min-order-bound}
 v(s)\leq\A(v)\mult_x(s).
\end{equation}
In particular, for $s\in R_m$ one has
$0\leq v(s)/\A(v)\leq Cm$, with $C$ as in
Lemma~\ref{lem:min-initials}.
\end{lemma}

\begin{proof}
For an effective divisor with local equation $g$ of multiplicity $r>0$
at a smooth point, the elementary bound
$\lct_x(g)\geq1/r$ suffices. One analytic proof uses Weierstrass
preparation in a general coordinate direction. On a smaller
polydisc, $g$ is a unit times a monic polynomial of degree $r$ in that
coordinate. For any $c<1/r$, factor this polynomial into its $r$
complex roots for each fixed choice of the other coordinates. H\"older's
inequality bounds the integral of its absolute value to the power
$-2c$ by the product of integrals with exponents $-2cr$. These
one-variable integrals are uniformly finite because $cr<1$ and the
roots remain bounded. Fubini's theorem proves local integrability.
The unit case has infinite threshold. Computing the threshold on a
log resolution now gives \eqref{eq:min-order-bound} for divisorial
orders with centers containing $x$, and rational approximation gives
it for the admitted monomial tests.
\end{proof}

\begin{lemma}[Uniform sublevel bounds]
\label{lem:min-coercivity}
Suppose that $P$ is not generated at $x$. There are numbers
$0<t_0<n+1$, $c>0$, $0<\delta<1$, $0<a_*\leq a^*<\infty$, and
an integer $m_0$ such that, for every $m\geq m_0$ and every
$t\in[t_0,n+1]$,
\begin{equation}
\label{eq:min-uniform-bounds}
 c\leq f_{m,t}(x)\leq1-\delta.
\end{equation}
Moreover, any pair in \eqref{eq:min-objective} centered with image
containing $x$ and satisfying
$F_{m,t}(v,\mathbf s)\leq1-\delta/2$ has
\begin{equation}
\label{eq:min-discrepancy-bounds}
 a_*\leq\A(v)\leq a^*.
\end{equation}
All these constants are independent of $m$ and $t$ in the indicated
ranges.
\end{lemma}

\begin{proof}
By Proposition~\ref{prop:ordinary-gap}, there is $\epsilon>0$ such
that, for all sufficiently large $m$,
\[
 \frac1{mN_m}\sum_{\gamma\in H_m}|\gamma|_1
       \geq\frac n{n+1}+3\epsilon.
\]
Choose $t_0<n+1$ sufficiently close that
$n/t_0\leq n/(n+1)+\epsilon$. Take a basis with distinct ordinary
initials, and use $v=\lambda\ord_x$, where $\ord_x$ is the
ordinary point order and $\A(\ord_x)=n$. In dimension one this is
the order of the point itself. The linear coefficient at $\lambda=0$
in its objective is at most $-2\epsilon$, uniformly in $m$ and $t$.
By \eqref{eq:min-multiplicity}, the quantities
$n/t-\mult_x(s_i)/m$ are uniformly bounded, so the quadratic remainder
in their exponential expansion is uniformly $O(\lambda^2)$. A fixed
sufficiently small $\lambda>0$ therefore gives
$F_{m,t}\leq1-\delta$ for a fixed $\delta>0$.

We next give a uniform supply of low-order jets. Choose $k>0$ such that
$kL$ is very ample, and sections $s_0,s_1,\ldots,s_n$ of $kL$ such
that $s_0(x)\ne0$ and $z_i=s_i/s_0$ are regular parameters at $x$.
For each residue class modulo $k$, choose a fixed sufficiently large
integer $b_r$ in that class for which $b_rL$ is globally generated,
and choose $\tau_r\in H^0(X,b_rL)$ nonzero at $x$. Write
$m=qk+b_r$ using the corresponding residue. For $|\alpha|_1\leq q$,
the sections
\[
 \tau_r s_0^{q-|\alpha|_1}s_1^{\alpha_1}\cdots s_n^{\alpha_n}
\]
have germs equal to a common unit times $z^\alpha$ in a local frame
of $mL$. In particular, if
\[
 r_m=\min\left(q,\left\lfloor\frac{m}{2(n+1)}\right\rfloor\right),
\]
their jets for $|\alpha|_1\leq r_m$ are linearly independent modulo
ordinary order greater than $m/(2t)$. For large $m$, $r_m$ is at least
a fixed positive multiple of $m$. The Hilbert asymptotic consequently
gives a constant $0<\kappa<1$, independent of $m,t$, with
\begin{equation}
\label{eq:min-jet-codimension}
 \codim_{R_m}\{s:\mult_x(s)>m/(2t)\}\geq\kappa N_m.
\end{equation}

Put $a=\A(v)$. By Lemma~\ref{lem:min-order-bound}, the subspace
$\{s:v(s)>ma/(2t)\}$ is contained in the subspace in
\eqref{eq:min-jet-codimension}. Thus at least $\kappa N_m$ vectors
of any basis, whether or not it is adapted to $v$, have
$v(s_i)\leq ma/(2t)$. It follows that
\begin{equation}
\label{eq:min-coercive-exponential}
 F_{m,t}(v,\mathbf s)
   \geq\kappa\exp\!\left(\frac a{2t}\right)
   \geq\kappa\exp\!\left(\frac a{2(n+1)}\right).
\end{equation}
This proves the lower bound in \eqref{eq:min-uniform-bounds}, with
$c=\kappa$, and bounds $a$ above on the stated sublevel set.

Finally, \eqref{eq:min-order-bound} and
\eqref{eq:min-multiplicity} show that
$F_{m,t}(v,\mathbf s)\geq e^{-Ca}$. Choose $a_*>0$ sufficiently
small that $e^{-Ca_*}>1-\delta/2$. No pair in the sublevel set can
have $a<a_*$. Enlarging the upper bound if necessary gives
$a^*\geq a_*$ and proves all the assertions.
\end{proof}

\begin{lemma}[Closed realizability loci]
\label{lem:min-flags}
Fix $m$ with $N_m>0$ and a decreasing list
$\alpha_1\geq\cdots\geq\alpha_{N_m}\geq0$ of rational numbers.
The set of points $z\in X$ for which some ordered basis and some
prime-divisor order $v$, with $z\in c_X(v)$, satisfy
\begin{equation}
\label{eq:min-rational-list}
 \frac{v(s_i)}{\A(v)}\geq\alpha_i
       \qquad(1\leq i\leq N_m)
\end{equation}
is Zariski closed. There is also a closed incidence locus in the
product of $X$ with the complete flag variety of $R_m$ recording
realizability by bases adapted to a given flag.
\end{lemma}

\begin{proof}
Let $\mathcal F$ be that projective flag variety. A flag has subspaces
$U_i\subset R_m$ of dimensions $i$ and associated base ideals
$\mathfrak b_i\subset\cO_X$, obtained by evaluation after twisting
by $-mL$. If \eqref{eq:min-rational-list} holds, the flag spanned by
the first $i$ basis vectors satisfies
\begin{equation}
\label{eq:min-flag-inequalities}
 v(\mathfrak b_i)\geq\alpha_i\A(v)
       \qquad(1\leq i\leq N_m),
\end{equation}
since the list is decreasing. Conversely these inequalities imply
\eqref{eq:min-rational-list} for every basis adapted to the flag.

If at least one $\alpha_i$ is positive, choose a positive integer $M$
such that $M/\alpha_i$ is integral for every positive entry, and set
\[
 \mathfrak a_\alpha
    =\sum_{i:\alpha_i>0}\mathfrak b_i^{\,M/\alpha_i}.
\]
The simultaneous inequalities \eqref{eq:min-flag-inequalities} are
equivalent to $v(\mathfrak a_\alpha)\geq M\A(v)$. A divisorial
witness with center containing $z$ exists if and only if
\begin{equation}
\label{eq:min-flag-threshold}
 \lct_z(X,\mathfrak a_\alpha)\leq1/M,
\end{equation}
because the local threshold is computed by a component on a log
resolution. The universal subbundles on $\mathcal F$ give algebraic
families of generators for the base ideals on $X\times\mathcal F$,
and hence for their indicated powers and sum. Lemma~\ref{lem:lct-family}
applied to \eqref{eq:min-flag-threshold} gives the closed incidence
locus. Its projection to $X$ is closed by properness of $\mathcal F$.
If all entries vanish, both loci are the whole parameter spaces.
\end{proof}

\begin{proposition}[Existence and uniform bounds for minimizers]
\label{prop:minimizers}
Suppose that $P$ is not generated at $x$. For the constants and ranges
in Lemma~\ref{lem:min-coercivity}, $f_{m,t}(x)$ satisfies
\eqref{eq:min-uniform-bounds} and is attained by an ordered basis
of $R_m$ and a nonzero real-weight monomial test on a projective log
resolution of the divisors of that basis. Every minimizing test has
\[
 a_*\leq\A(v)\leq a^*.
\]
More generally the same bounds hold eventually along every minimizing
sequence. In particular, the positive lower bound $a_*$ is independent
of $m\geq m_0$ and $t\in[t_0,n+1]$.
\end{proposition}

\begin{proof}
Only attainment remains to be proved. Fix $m,t$ in the indicated
ranges and take a sequence of divisorial pairs whose objectives
decrease to $f_{m,t}(x)$. Eventually they lie in the sublevel set of
Lemma~\ref{lem:min-coercivity}. Order each basis by decreasing
normalized values and pass to a subsequence for which
\[
 a_b=\A(v_b)\longrightarrow a\in[a_*,a^*],\qquad
 q_{b,i}=\frac{v_b(s_{b,i})}{\A(v_b)}\longrightarrow q_i.
\]
This is possible because $0\leq q_{b,i}\leq Cm$. The limit list is
decreasing and
\begin{equation}
\label{eq:min-limit-objective}
 f_{m,t}(x)=\frac1{N_m}\sum_i
       \exp\!\left(\frac at-\frac{aq_i}m\right).
\end{equation}
At least one $q_i$ is positive: otherwise the right side would be
$e^{a/t}>1$, contrary to \eqref{eq:min-uniform-bounds}.

Choose decreasing nonnegative rational lists
$\alpha^{(k)}=(\alpha_i^{(k)})$ which increase componentwise to
$(q_i)$ and are strictly below $q_i$ whenever $q_i>0$. Each list is
satisfied by all sufficiently late pairs in the minimizing sequence.
By the incidence assertion of Lemma~\ref{lem:min-flags}, the flags
admitting a divisorial witness at $x$ for that list form a nonempty
closed subset $\mathcal F_k$ of $\mathcal F$. The subsets are nested,
so projectivity gives a flag in their intersection. Fix a basis
$s_1,\ldots,s_{N_m}$ adapted to this one flag.

Resolve the divisors of these sections on a single smooth projective
model $Y\longrightarrow X$, including the exceptional divisors in
the SNC support. For each sufficiently large $k$, choose a divisorial
witness for the fixed flag and the list $\alpha^{(k)}$, and exhibit it
on a common smooth model dominating $Y$. Its retraction
to this fixed SNC model, as in Lemma~\ref{lem:retraction}, preserves
all the section orders, cannot increase discrepancy, and has center
on $X$ containing $x$. The retraction is nonzero because the list
has a positive entry. Normalize it to discrepancy one. It then
satisfies
\begin{equation}
\label{eq:min-retracted-lists}
 \widetilde v_k(s_i)\geq\alpha_i^{(k)}
       \qquad(1\leq i\leq N_m).
\end{equation}

There are finitely many SNC strata and their components on $Y$.
After passing to a subsequence, the retractions use the same stratum
component, with coordinate divisors indexed by a fixed set $J$.
Their normalized weights belong to the compact simplex
\[
 \left\{(u_j)_{j\in J}:u_j\geq0,\quad
                \sum_{j\in J}\A(E_j)u_j=1\right\};
\]
all coefficients $\A(E_j)$ are strictly positive because $X$ is
smooth. Pass to a limit of the weights. A weight may become zero.
We take the component of the remaining positive-support intersection
that contains the old stratum component. Its closure contains the old
center, so its image still contains $x$.
The limit is a nonzero monomial test $\widetilde v$
of discrepancy one. On the fixed resolution, the section orders are
linear in these weights. Consequently
\eqref{eq:min-retracted-lists} gives
$\widetilde v(s_i)\geq q_i$ for every $i$.

Set $v=a\widetilde v$. Comparing with
\eqref{eq:min-limit-objective} yields
\[
 F_{m,t}(v,\mathbf s)
 \leq\frac1{N_m}\sum_i
       \exp\!\left(\frac at-\frac{aq_i}m\right)
 =f_{m,t}(x).
\]
The reverse inequality holds by definition of the infimum, so this
pair attains it. The claimed discrepancy bounds for all minimizers
and eventually for minimizing sequences follow directly from the
uniform sublevel bounds.
\end{proof}

The same closed realizability loci also compare the infimum along a
possible center. This comparison will allow the tangential variation in
Section~\ref{sec:centers} to be tested at a very general point.

\begin{lemma}[Specialization comparison]
\label{lem:specialization}
For fixed $m$ with $N_m>0$, let $W\subset X$ be a closed irreducible
subvariety containing $x$. At a very general closed point $y\in W$,
one has, for every $t>0$,
\begin{equation}
\label{eq:min-specialization}
 f_{m,t}(y)\geq f_{m,t}(x).
\end{equation}
\end{lemma}

\begin{proof}
There are countably many rational lists of the kind in
Lemma~\ref{lem:min-flags}. For each, intersect its closed realizability
locus with $W$, and exclude that intersection if it is proper in $W$.
The complement contains very general closed points over $\C$.
For any remaining $y$, every such list realized at $y$ is realized
at $x$.

Take a divisorial test and a basis at $y$, reorder the basis so that
$q_i=v(s_i)/\A(v)$ is decreasing, and put $a=\A(v)>0$. Approximate
$(q_i)$ componentwise from below by decreasing rational lists
$(\alpha_i^{(b)})$. For each list, choose a realizing pair at $x$ and
scale its test to discrepancy exactly $a$. Its objective is at most
\[
 \frac1{N_m}\sum_i
       \exp\!\left(\frac at-\frac{a\alpha_i^{(b)}}m\right).
\]
Taking the limit gives $f_{m,t}(x)\leq F_{m,t}(v,\mathbf s)$.
Infimizing over the divisorial pairs at $y$ proves
\eqref{eq:min-specialization}. The excluded loci do not depend on
$t$, which explains the simultaneous statement for all positive $t$.
\end{proof}
